\subsection{Sprint 4}

Por se tratar da sprint final, não havia muitas tarefas novas, mas sim débitos
técnicos de sprints passadas, além de complicações com a integração que deveriam
ser resolvidas. Da minha parte tinha a página de administração, na qual eu
precisava atualizar parâmetros e outros atributos do display conforme redefinições
de regra de negócio e alterações do backend, além de adicionar qualquer botão que
faltasse na tela de administração.

O time foi capaz de entregar praticamente tudo, incluindo tarefas novas que
iam sendo criadas conforme reflexão do próprio time relativas à experiência de
usuário. Da minha parte fiquei um bom tempo esperando a rota do backend ser
implementada para a parte de administração, já que era uma das páginas mais
complexas da parte do backend.

Enquanto esperava o backend finalizar a rota, me prontifiquei para realizar pequenas mudanças de qualidade de vida para a página,
como adicionar um botão de relatório integrado com o backend, reunir com outros
colegas para adicionar o drag and drop e a criação de usuários por arquivo e outras
tarefas do tipo. Por fim participei frequentemente das reuniões de alinhamento do
time fora da aula para garantir que me manteria informado e sempre ajudando o time nas tarefas necessárias.

Más práticas por parte do backend como, não padronizar o nome esperado do
payload e deixar a documentação do swagger mal-feita, foram os principais
problemas encontrados nesta última sprint, já que para realizar as integrações era necessário
entender o código em NestJS \cite{nestjs-docs} ou perguntar detalhadamente para o integrante
responsável pelo código e isso foi responsável por diminuir um pouco a velocidade
com que o desenvolvimento ocorreu.

Da minha parte acredito que poderia ter ido atrás de aprender o mínimo de NestJS
\cite{nestjs-docs},
já que isso teria facilitado muito a legibilidade do código do backend e, portanto,
facilitado a parte de integração. Ao mesmo tempo pensando pelos olhos de quem
faz o backend, que é onde planejo me concentrar mais na próxima Ages, entendi
muito bem qual a importância de uma boa documentação e um Swagger fácil de
entender, já que percebi na prática o quanto ele fez falta para alguém do frontend
que não faz ideia de como funciona a linguagem do backend.

Com a última sprint desta Ages finalizada, só basta pensar na próxima, como
citado acima, espero me envolver muito mais no backend para chegar de Ages 3
com uma boa experiência nos dois lados. Além disso acredito que a função de Ages 2 de
cuidar do banco de dados combina muito bem com o desenvolvimento backend, já
que os dois referem-se bem restritamente à regra de negócio. Além disso estou
querendo aprender Java Spring \cite{spring-docs} por fora da aula, por interesse próprio, o que me
daria um bom conhecimento básico do backend.
